\chapter{متى نستكشف ومتى نقرر: نضيف \nt{NORA}}
\chaptermark{نضيف \nt{NORA}}
\label{sec:nora}
\begin{chaptergoals}
\item كيف ينقل مقبض كسب واحد العصبون من مضخّم تماثلي إلى مقارِن؛
\item لماذا لا يفيد رفع الكسب إلا ضد الضجيج الذي ينشأ بعد المضخّم؛
\item كيف ينقل الكسب شبكة الاختيار من التروّي إلى القرار دون أن يمسّ أي وزن؛
\item لماذا الكسب مقلوب درجة الحرارة، ولماذا تتبع المكافأة حرف U مقلوبًا بدلالة الكسب.
\end{chaptergoals}

\section{ما الذي ينقص}

في الفصل~\ref{sec:dofa} تعلّمت الشبكة بفضل الضجيج: كانت الانحرافات العشوائية في النشاط بحثًا، وكان \nt{DOPA} يخبر هل صار الوضع بها أفضل. لكن بقي هناك معامل واحد غير محدد: \emph{كم} نبحث. إذا كان الضجيج قويًا جدًا تخبطت الشبكة ولم تنتفع بما تعرفه. وإذا كان ضعيفًا جدًا اختارت الشيء نفسه مرة بعد مرة، ولم تلاحظ أن خيارًا قريبًا صار أفضل. في التعلّم المعزَّز يسمى هذا الاختيارَ بين الاستغلال والاستكشاف، ولكل خوارزمية مقبض له. وفي الفصل~\ref{sec:neurontypes} افترضنا أن \nt{NORA} هو الذي يدير هذا المقبض في الدماغ.

عصبونات النورأدرينالين عند الإنسان بضع عشرات الآلاف (الجدول~\ref{tab:types})، وكلها تقريبًا مجموعة في تجمّع صغير واحد، أما مخارجها فتتفرق في الدماغ كله تقريبًا، وإن تبيّن أن الأجزاء المختلفة من هذا التجمع تخدم جزئيًا مناطق مختلفة~\cite{Poe2020}. هذه قناة بثّ أضيق حتى من \nt{DOPA}. وهي تعمل في وضعين~\cite{AstonJones2005}. في الوضع \emph{الخلفي} تطلق العصبونات النبضات بانتظام، بتردد من أجزاء النبضة إلى بضع نبضات في الثانية، ويتغير هذا التردد ببطء مع درجة اليقظة. وفي الوضع \emph{الطوري} تستجيب برشقة قصيرة لحدث مهم للمهمة الجارية، نحو لحظة اتخاذ القرار. ونقطة الاختبار المريحة لكن غير الدقيقة هي قطر الحدقة: فهو يتغير مع نشاط هذه العصبونات، لكنه يتغير أيضًا مع نشاط مناطق أخرى~\cite{Joshi2016} ومع المخارج الكولينية في القشرة~\cite{Reimer2016}. الحدقة تبيّن المستوى العام للإثارة، لا \nt{NORA} وحده.

\section{معامل الكسب}

المقيس هو التالي: النورأدرينالين يكبت النشاط الخلفي للعصبونات أكثر مما يكبت استجابتها للمنبّه، فترتفع نسبة الإشارة إلى الخلفية~\cite{Foote1975NE}. أما أن انحدار خاصية العصبون المفرد يُضرب حرفيًا في معامل، فلا يلزم من التجربة. وكما في الفصل~\ref{sec:neurontypes}، نأخذ نموذجًا بسيطًا ينقل هذا الأثر: النورأدرينالين يغيّر انحدار خاصية النقل عند المتلقي~\cite{ServanSchreiber1990}. عندئذ تعطي المداخل الضعيفة تحت العتبة (الخلفية) أقل، والقوية أكثر. ليستجب العصبون بتردد
\begin{empheq}[box=\keybox]{equation}
r \;=\; F\bigl(g\,(u-\theta)\bigr), \qquad F(z)=\frac{r_{\max}}{1+e^{-z}} ,
\label{eq:gain}
\end{empheq}
حيث $u$ تيار الدخل، و$\theta$ العتبة، و$g$ معامل الكسب الذي يتحكم فيه \nt{NORA} في النموذج. عند $g$ صغير يتبع التردد الدخلَ بسلاسة في مدى واسع: العصبون \emph{مضخّم تماثلي}. وعند $g$ كبير يعطي أي دخل فوق العتبة تقريبًا التردد الكامل، ويعطي ما تحتها صفرًا: العصبون \emph{مقارِن}، أي عنصر شبه رقمي (الشكل~\ref{fig:gaincurves}). مقبض واحد ينقل العصبون نفسه بين وضعين تنفذهما في الإلكترونيات دارات مختلفة.

\begin{figure}[t]
\centering
\begin{tikzpicture}[>=Stealth,x=1.1cm,y=2.4cm]
\draw[->,gray!70!black] (-3.3,0) -- (3.4,0) node[below left,black] {\small $u$};
\draw[->,gray!70!black] (-3.3,0) -- (-3.3,1.2) node[left,black] {\small $r$};
\draw[densely dashed,gray!60!black] (0,0) -- (0,1.1);
\node[below,font=\small] at (0,0) {$\theta$};
\draw[densely dotted,gray!60!black] (-3.3,1) -- (3.2,1) node[right,black,font=\small] {$r_{\max}$};
\draw[very thick,blue!60!black] plot[domain=-3.2:3.2,samples=60] (\x,{1/(1+exp(-0.8*\x))});
\draw[very thick,orange!85!black] plot[domain=-3.2:3.2,samples=60] (\x,{1/(1+exp(-2*\x))});
\draw[very thick,red!70!black] plot[domain=-3.2:3.2,samples=120] (\x,{1/(1+exp(min(-8*\x,9)))});
\node[blue!60!black,font=\small,anchor=west] at (0.5,0.12) {$g$ صغير: مضخّم};
\node[red!70!black,font=\small,anchor=east] at (-0.3,0.85) {$g$ كبير: مقارِن};
\end{tikzpicture}
\caption{خاصية النقل~\eqref{eq:gain} عند ثلاث قيم لمعامل الكسب $g$.}
\label{fig:gaincurves}
\end{figure}

هل يساعد الكسب الكبير على مقاومة الضجيج؟ ليس دائمًا. ليختلف دخلان بمقدار $\Delta u$، والمطلوب أن نقول أيهما أكبر. إذا أُضيف الضجيج $\sigma_{\text{pre}}$ إلى الدخل \emph{قبل} المضخّم، تضخّمت الإشارة والضجيج معًا، ولم تتغير نسبتهما. أما إذا أُضيف الضجيج $\sigma_{\text{post}}$ \emph{بعد} المضخّم، من المشابك غير الموثوقة ومن النبضات العشوائية للشبكة نفسها، كما في الفصل~\ref{sec:code}، فإن نسبة الإشارة إلى الضجيج
\begin{empheq}[box=\keybox]{equation}
d' \;=\; \frac{g\,\Delta u}{\sqrt{g^2\sigma_{\text{pre}}^2+\sigma_{\text{post}}^2}}
\label{eq:gainsnr}
\end{empheq}
تزداد مع $g$ إلى أن يتساوى $g\sigma_{\text{pre}}$ مع $\sigma_{\text{post}}$~\cite{ServanSchreiber1990}.

\begin{keyblock}
\begin{proposition}[متى يساعد الكسب]
\label{prop:gainnoise}
حين يرفع \nt{NORA} معامل الكسب، فإنه يحسّن التمييز بين المداخل ضد الضجيج الذي ينشأ \emph{بعد} المضخّم، في الشبكة نفسها، فقط. أما الضجيج الذي جاء مع الدخل فلا يزيله الكسب.
\end{proposition}
\end{keyblock}

\section{الكسب في الاختيار بين اثنين}

لنعد إلى الاختيار من الفصل~\ref{sec:gaba}: مجموعتان من عصبونات \nt{GLUT} مدخلاهما $I_1$ و$I_2$ تثبّط إحداهما الأخرى بقوة $\beta$. الآن لكل مجموعة معامل كسب $g$: في معادلات الاختيار~\eqref{eq:wta} يُضرب ما بين القوسين في $g$. ما دامت المجموعتان تعملان، تُستخرج الترددات المستقرة من $r_1=g(I_1-\beta r_2)$ و$r_2=g(I_2-\beta r_1)$. بجمع هاتين المساواتين وطرحهما نحصل على المجموع $r_1+r_2=g(I_1+I_2)/(1+g\beta)$ والفرق $r_1-r_2=g(I_1-I_2)/(1-g\beta)$. ونسبتهما تبيّن مدى حدة تفضيل الشبكة لدخل على آخر:
\begin{empheq}[box=\keybox]{equation}
\frac{r_1-r_2}{r_1+r_2} \;=\; \varepsilon\,\frac{1+g\beta}{1-g\beta}, \qquad \varepsilon=\frac{I_1-I_2}{I_1+I_2}, \quad g\beta<1 .
\label{eq:wtagain}
\end{empheq}
الفرق الصغير بين المدخلين $\varepsilon$ يتضخم $(1+g\beta)/(1-g\beta)$ مرة، وهذا العامل يكبر بلا حد حين يقترب $g\beta$ من الواحد. وعند $g\beta>1$ تكون الحالة التي تعمل فيها المجموعتان غير مستقرة، كما في القضية~\ref{prop:wta}: تفوز واحدة وتصمت الأخرى (الشكل~\ref{fig:pitchfork}).

\begin{figure}[t]
\centering
\begin{tikzpicture}[>=Stealth,x=3.2cm,y=1.5cm]
\draw[->,gray!70!black] (0,0) -- (2.15,0) node[below left,black] {\small $g\beta$};
\draw[->,gray!70!black] (0,-1.25) -- (0,1.35) node[above,black] {\small $(r_1-r_2)/(r_1+r_2)$};
\draw[gray!60!black] (1,-0.04) -- (1,0.04); \node[below right,font=\small] at (1,0) {$1$};
\node[left,font=\scriptsize] at (0,1) {$1$}; \node[left,font=\scriptsize] at (0,-1) {$-1$};
\draw[very thick,blue!60!black] plot[domain=0:0.905,samples=60] (\x,{0.05*(1+\x)/(1-\x)}) -- (1,1) -- (2.05,1);
\draw[very thick,blue!60!black] (1.105,-1) -- (2.05,-1);
\draw[thick,densely dashed,gray!60!black] (1,0) -- (2.05,0);
\node[font=\small,align=center] at (0.45,-0.62) {تتروّى:\\المجموعتان\\تعملان};
\node[font=\small,align=center] at (1.55,0.45) {قرّرت:\\تعمل واحدة};
\node[font=\small,blue!60!black,anchor=south west] at (1.05,-0.97) {فاز الدخل الأضعف أولًا، فيبقى فائزًا};
\end{tikzpicture}
\caption{الاختيار بين اثنين عند معاملات كسب مختلفة، والمدخلان يختلفان بـ$\varepsilon=0.05$. عند $g\beta>1$ يكون القراران كلاهما مستقرين (الخطان المتصلان؛ وفوز الدخل الأضعف مستقر عند $g\beta>(1+\varepsilon)/(1-\varepsilon)\approx1.1$)، وتتمسك الشبكة بما اتخذته منهما؛ أما الحالة المتناظرة (الخط المتقطع) فغير مستقرة.}
\label{fig:pitchfork}
\end{figure}

\begin{keyblock}
\begin{proposition}[الكسب يبدّل وضع الاختيار]
\label{prop:gainwta}
في شبكة اختيار بتثبيط متبادل $\beta$ ينقل معامل الكسب $g$ الشبكة عبر الحد $g\beta=1$. تحته ”تتروّى“ الشبكة: تعمل المجموعتان، ويتضخم فرق المدخلين وفق~\eqref{eq:wtagain}. وفوقه ”قرّرت“ الشبكة: تعمل مجموعة واحدة، ويثبت القرار. وبتغيير $g$ يستطيع \nt{NORA} أن ينقل الشبكة بين هذين الوضعين دون أن يمسّ أي وزن.
\end{proposition}
\end{keyblock}

\section{الكسب بوصفه درجة حرارة}

لنضف الآن إلى الاختيار ضجيجًا ينشأ في الشبكة نفسها، بعد المضخّم. أي المجموعتين تفوز تحدده إشارة المقدار $g(u_A-u_B)+\eta$، حيث $u_A-u_B$ فرق المدخلين، أي فرق الأوزان المتعلَّمة من الفصل~\ref{sec:dofa}، و$\eta$ ضجيج مقياسه $\sigma$. إذا كان الضجيج موزعًا وفق القانون اللوجستي (والمنحنى يكاد يكون نفسه للتوزيع الغاوسي)، فإن احتمال اختيار $A$ هو
\begin{empheq}[box=\keybox]{equation}
P(A) \;=\; \frac{1}{1+e^{-g\,(u_A-u_B)/\sigma}} .
\label{eq:softmax}
\end{empheq}
سيعرف المبرمج هنا الاختيار بدالة \foreignlanguage{english}{softmax} بدرجة حرارة $T=\sigma/g$. معامل الكسب هو مقلوب درجة الحرارة. عند $g$ صغير لا يُختار حتى الخيار الأفضل بوضوح أكثر من الأسوأ إلا قليلًا: الشبكة تستكشف كثيرًا. وعند $g$ كبير يُختار أفضل خيار معروف دائمًا تقريبًا: الشبكة تستغل ما تعرفه.

لنوضح ما هو $g$ في~\eqref{eq:wtagain} و~\eqref{eq:softmax}: إنه الكسب \emph{الفعّال} لفرق مدخلات المهمة الجارية لحظة الاختيار. ووضعا \nt{NORA} يؤثران فيه تأثيرين متعاكسين. الرشقة الطورية لحظة القرار ترفعه، فتثبّت الشبكة اختيارها. أما النشاط الخلفي المرتفع فيرافق الاستكشاف على العكس: عند البشر تكون الحدقة القاعدية أوسع قبل الاختيار العشوائي~\cite{JepmaNieuwenhuis2011}، وعند الجرذان ينقل دخل \nt{NORA} إلى القشرة الحزامية الأمامية الاختيارَ إلى وضع عشوائي~\cite{Tervo2014stochastic}. إذن يقابل \nt{NORA} الخلفيَّ المرتفع $g$ فعّال \emph{صغير}، ويقابل الرشقةَ $g$ كبير.

\section{كم نستكشف}

أي كسب أفضل؟ تحققنا من ذلك على نموذج. تختار الشبكة أحد فعلين وفق~\eqref{eq:softmax}؛ وكل فعل يجلب مكافأة باحتمال ما، والشبكة تقدّره من مكافآت الفعل المختار، كما في الفصل~\ref{sec:dofa}. ومن حين لآخر يتغير العالم: يُسحب الاحتمالان من جديد. قد يتغير الفعل المختار، وقد يتغير الفعل الذي لم تجربه الشبكة منذ زمن؛ ولا تعرف الشبكة بالثاني إلا بالاستكشاف. يبيّن الشكل~\ref{fig:invertedU} نسبة أفضل مكافأة ممكنة التي تنالها الشبكة عند قيم ثابتة مختلفة لـ$g$.

\begin{figure}[t]
\centering
\begin{tikzpicture}[>=Stealth,x=1.2cm,y=1cm]
\draw[->,gray!70!black] (0,0) -- (7.6,0) node[below left,black] {\small $g$};
\draw[->,gray!70!black] (0,0) -- (0,4.4) node[above,black,font=\small] {نسبة أفضل مكافأة};
\foreach \x/\l in {0/0.5,1/1,2/2,3/4,4/8,5/16,6/32,7/64}{\draw[gray!60!black] (\x,-0.06) -- (\x,0.06); \node[below,font=\scriptsize] at (\x,0) {\l};}
\foreach \v/\y in {0.8/0.8,0.9/2.4,1.0/4.0}{\draw[gray!60!black] (-0.06,\y) -- (0.06,\y); \node[left,font=\scriptsize] at (0,\y) {\v};}
\draw[very thick,blue!60!black] plot[mark=*,mark size=1.3pt] coordinates {(0,0.368) (1,0.880) (2,1.728) (3,2.784) (4,3.456) (5,3.616) (6,3.488) (7,3.264)};
\draw[very thick,red!70!black] plot[mark=*,mark size=1.3pt] coordinates {(0,0.368) (1,0.672) (2,1.184) (3,1.840) (4,2.176) (5,1.920) (6,1.456) (7,1.104)};
\node[blue!60!black,font=\small,anchor=south] at (5,3.7) {عالم هادئ};
\node[red!70!black,font=\small,anchor=north] at (5.1,1.25) {عالم سريع التغير};
\draw[->,thick,blue!60!black] (5,3.25) -- (5,3.55);
\draw[->,thick,red!70!black] (4,1.8) -- (4,2.1);
\node[font=\small\itshape,gray!35!black,anchor=west] at (0.15,2.3) {تبحث عشوائيًا};
\node[font=\small\itshape,gray!35!black] at (6.45,2.55) {لا تلاحظ التغيرات};
\end{tikzpicture}
\caption{نسبة أفضل مكافأة ممكنة عند معامل كسب ثابت $g$ (المقياس على محور $g$ لوغاريتمي). المنحنى الأزرق: يتغير العالم في المتوسط مرة كل ألف محاولة، والأحمر: مرة كل عشرين. تشير السهام إلى أفضل $g$: في العالم السريع التغير هو أصغر بمرتين. المتوسط على 60 محاكاة في كل منها 4000 محاولة.}
\label{fig:invertedU}
\end{figure}

عند $g=0.5$ تكاد الشبكة لا تستعمل معرفتها، فتنال نحو ثلاثة أرباع الممكن، وعند $g$ كبير جدًا تفوتها التغيرات. أفضل قيمة: $g=16$ حين يتغير العالم مرة كل ألف محاولة (النسبة $0.976$)، و$g=8$ حين يتغير مرة كل عشرين ($0.886$).

\begin{keyblock}
\begin{proposition}[حرف U مقلوب]
\label{prop:explore}
المكافأة التي تجمعها الشبكة تعتمد على معامل الكسب اعتمادًا على شكل حرف U مقلوب: $g$ الصغير جدًا يهدر المحاولات في البحث، والكبير جدًا لا يلاحظ التغيرات. وأفضل $g$ يصغر كلما تغيّر العالم أسرع.
\end{proposition}
\end{keyblock}

ويُلاحَظ حرف U المقلوب في التجارب أيضًا: تحل الحيوانات المهمة على أفضل وجه عند نشاط خلفي معتدل لعصبونات \nt{NORA} مع استجابات طورية واضحة، أما عند النشاط الخلفي المرتفع فتتشتت وتتنقل بين المهام~\cite{AstonJones2005}. وهذا يتفق مع الفرق بين الوضعين في القسم السابق: الرشقة الطورية لحظة القرار تعطي $g$ فعالًا كبيرًا في الوقت الذي يجب فيه الاختيار تحديدًا، أما النشاط الخلفي المرتفع بلا رشقات فيضخّم كل شيء دون تمييز، بما في ذلك المداخل الدخيلة، فيهبط $g$ الفعّال للمهمة الجارية، وهذا هو الاستكشاف.

\section{من يدير المقبض}

ما دام أفضل كسب يعتمد على سرعة تغيّر العالم، فمن الحسن أن يُضبط. لكن من أين تعرف عصبونات \nt{NORA} أن العالم تغيّر؟ العلامة الطبيعية هي \emph{المفاجأة}: صارت أخطاء التنبؤ أكبر من المعتاد. من المهم التمييز بين نوعين من عدم اليقين. إذا كانت المكافأة عشوائية ببساطة، فالأخطاء كبيرة دائمًا، وهذا متوقع. أما إذا ازدادت الأخطاء فجأة مقارنة بحجمها المعتاد، فقد تغيّر شيء ما. ووفق إحدى الفرضيات، يخبر \nt{NORA} بعدم اليقين غير المتوقع هذا تحديدًا، أما المتوقع فيخبر به \nt{ACET}~\cite{YuDayan2005}.

ماذا نفعل بهذه الإشارة؟ يمكن خفض الكسب والاستكشاف أكثر. ويمكن رفع معدل التعلّم لنسيان التقديرات القديمة أسرع: تبيّن التجارب أن الناس يتعلمون أسرع بعد تغيّر حاد، وأن الحدقة تتسع عندئذ~\cite{Nassar2012}؛ ونذكّر بأن الحدقة مؤشر لا على \nt{NORA} وحده، بل على \nt{ACET} أيضًا~\cite{Reimer2016}. ويمكن فعل الأمرين معًا: وفق فرضية أخرى تعمل رشقة \nt{NORA} الطورية \emph{مقاطعةً}، أي إشارة تعيد بها الشبكة ضبط حالتها الجارية وتتكيف مع الموقف الجديد~\cite{BouretSara2005}، كخط المقاطعة العام في المعالج.

لنقل بصراحة ما لم نجده. جربنا في النموذج طريقتي الضبط البسيطتين كلتيهما: خفض $g$ أو رفع معدل التعلّم حين يتجاوز الخطأ المنعَّم متوسطه الطويل الأمد. في عالم يثبت مدة طويلة تارة ويتغير كثيرًا تارة أخرى، أعطت أفضل إعدادات الطريقتين $0.926\text{--}0.927$ من أفضل مكافأة، أي ما يكاد يساوي أفضل $g$ ثابت ($0.925$ عند $g=8$) أو معدل تعلّم ثابتًا لكنه كبير بما يكفي ($0.928$)، وأعطت الإعدادات غير الموفقة أقل. المنحنيات في الشكل~\ref{fig:invertedU} مسطحة قرب القمة، وفي عالم كهذا لا يكاد يوجد ما يُضبط. ينبغي أن يؤتي الضبط ثماره حيث تتطلب الأوضاع المختلفة إعدادات مختلفة جدًا. أما قانون التحكم الذي ينفذه \nt{NORA} بالضبط فلم يُحدَّد بعد.

\section{الخلاصة}

\begin{table}[t]
\centering
\footnotesize
\setlength{\tabcolsep}{4pt}
\renewcommand{\arraystretch}{1.15}
\begin{tabular}{>{\raggedright}p{3.0cm}>{\raggedright}p{4.4cm}>{\raggedright\arraybackslash}p{6.0cm}}
\toprule
ما يفعله \nt{NORA} & كيف & ما هو معروف \\
\midrule
يغيّر انحدار استجابة العصبونات & معامل الكسب $g$ في~\eqref{eq:gain} & ارتفاع نسبة الإشارة إلى الضجيج مقيس؛ والنموذج الضربي تبسيط \\
يبدّل وضع الاختيار & ينقل الشبكة عبر $g\beta=1$ & ينتج من النموذج~\eqref{eq:wtagain}؛ ولا قياسات مباشرة للحد \\
يحدد كم نستكشف & $g$ مقلوب درجة الحرارة~\eqref{eq:softmax}؛ الخلفي $g$ فعّال صغير (استكشاف)، والرشقة كبير (قرار) & حرف U المقلوب ووضعا العمل مقيسة؛ والصلة بالاستكشاف فرضية قوية الأساس \\
يخبر بالتغيرات & عدم اليقين غير المتوقع & اتساع الحدقة وارتفاع معدل التعلّم بعد التغيرات مقيسان؛ وأنه إشارة \nt{NORA} فرضية \\
يقاطع ويعيد الضبط & رشقة طورية عند حدث غير متوقع & الرشقات عند الأحداث المهمة مقيسة؛ و”المقاطعة“ فرضية \\
\bottomrule
\end{tabular}
\caption{ما يضيفه \nt{NORA} إلى شبكة من \nt{GLUT} و\nt{GABA} و\nt{DOPA}.}
\label{tab:nora}
\end{table}

\nt{DOPA} يقول للشبكة \emph{إلى أين} تغيّر الأوزان. أما \nt{NORA} فلا يغيّر أي وزن، لكنه يقرر \emph{كيف} تستعمل الشبكة ما تعرفه (الجدول~\ref{tab:nora}): مقبض واحد، هو معامل الكسب، يحوّل العصبون من مضخّم إلى مقارِن، وشبكة الاختيار من ”متروّية“ إلى ”قررت“، والاختيار من استكشاف إلى استغلال. أما كيف يعرف \nt{NORA} سرعة تغيّر العالم فسؤال مفتوح.

\section{تمارين}

\begin{exercise}[\lvlA]
\label{ex:08:1}
\emph{مقبض واحد للدماغ كله.} (أ)~قدّر من الجدول~\ref{tab:types} كم عصبونًا في الدماغ يقابل عصبون \nt{NORA} واحدًا. (ب)~الضجيج قبل المضخّم $\sigma_{\text{pre}}=0.5$، وبعده $\sigma_{\text{post}}=4$. عند أي $g$ يبلغ $d'$ من~\eqref{eq:gainsnr} $90\,\%$ من حدّه؟
\end{exercise}

\begin{exercise}[\lvlB]
\label{ex:08:2}
\emph{الاختيار مع معامل كسب.} $\tau\,\dd r_1/\dd t=-r_1+g[I_1-\beta r_2]_+$، $\tau\,\dd r_2/\dd t=-r_2+g[I_2-\beta r_1]_+$، $\tau=20\ms$. (أ)~في كم من الزمن يخمد $r_1-r_2$ عند $g\beta=0.5$ و$0.9$؟ (ب)~عند $\varepsilon=0.05$ جد قيمة $g\beta$ التي تعمل تحتها المجموعتان، وتلك التي يستقر فوقها فوز الدخل الأضعف.
\end{exercise}

\begin{exercise}[\lvlA]
\label{ex:08:3}
\emph{\foreignlanguage{english}{Softmax} من الضجيج.} تتلقى كل مجموعة من $K$ مجموعة ضجيج غامبل خاصًا بها، $P(\eta_k<z)=\exp(-e^{-z/\sigma})$، ويفوز أكبر $gu_k+\eta_k$. جد $P(k)$. عند أي $g$ يُختار الأفضل من خيارين، مع $u_A-u_B=0.1$ و$\sigma=1$، في $90\,\%$ من الحالات؟
\end{exercise}

\begin{exercise}[\lvlB]
\label{ex:08:4}
\emph{تقدير أفضل كسب.} في نموذج الشكل~\ref{fig:invertedU} تكون احتمالات المكافآت بعد التغيّر موزعة بانتظام على $[0,1]$. تجرب الشبكة الفعل الأسوأ مرة كل $e^{g\Delta}$ محاولة؛ وبمساواة ذلك بـ$1/h$ نحصل على $g^*\approx\ln(1/h)/\bar\Delta$. جد $\bar\Delta=\langle|p_A-p_B|\rangle$ و$g^*$ عند $h=0.001$ و$0.01$ و$0.05$.
\end{exercise}

\begin{exercise}[\lvlB]
\label{ex:08:5}
\emph{نمذجة: الكسب وسرعة التغيرات.} جد $g^*(h)$ عند $h$ من $0.001$ إلى $0.2$ بنسخة من \texttt{models/\allowbreak nora\_explore.py} (وهي تكتب ملفات في المجلد الحالي). أين يصح تقدير التمرين~\ref{ex:08:4}، ولماذا يكفّ $g^*$ عن التناقص عند التغيرات السريعة؟
\end{exercise}

\begin{exercise}[\lvlB]
\label{ex:08:6}
\emph{تصميم: مقاطعة لشبكة الاختيار.} شبكة التمرين~\ref{ex:08:2} مع $\beta=2$ و$\tau=20\ms$ و$g=1$ تثبت على $r_1=0.9$، $r_2=0$، مع أن المدخلين صارا $I_1=0.9$ و$I_2=1.0$. يخفض \nt{NORA} مدة $T_p$ قيمة $g$ إلى $g_{\text{lo}}$. جد أصغر $T_p$ تحليليًا عند $g_{\text{lo}}=0$، وبالنمذجة عند $g_{\text{lo}}=0.25$.
\end{exercise}

\begin{exercise}[\lvlC]
\label{ex:08:7}
\emph{متى يؤتي الضبط ثماره.} يعرف عرّاف هل الحقبة هادئة أم متقلبة، ويضع في كل منها $g$ خاصًا بها. (أ)~جد بالنمذجة مكسبه على أفضل $g$ ثابت في عالم الفصل (حقب من $1000$ محاولة، $h=0.001$ و$0.05$). (ب)~ابنِ عالمًا يكون فيه المكسب أكبر، وجد الجزء الذي يستعيده منه الضبط بحسب المفاجأة.
\end{exercise}
